% GENERATED FILE. Edit canonical lessons, then run npm run generate:books.
% canonical-source-hash: fnv1a64:fcd9fc33b2c67c09
% canonical-lessons: TA-C184-panikkatti, TA-C184-kalai-unavu, TA-C184-matiya-unavu, TA-C184-iravu-unavu, TA-C184-kurippetu

\chapter{Ice, Breakfast, Lunch, Dinner, Notebook}
\label{ch:ta-ice-breakfast-lunch-dinner-notebook}
\hlchaptermodality{184}

\begin{chapteropening}
\textbf{By the end of this chapter:} \emph{I can say ice, breakfast, lunch, dinner and a notebook.}

Complete the last lesson of chapter 184: I can say ice, breakfast, lunch, dinner and a notebook.
\end{chapteropening}

\section[\texorpdfstring{\ta{பனிக்கட்டி} (paṉikkaṭṭi)}{பனிக்கட்டி (paṉikkaṭṭi)}]{\ta{பனிக்கட்டி} (paṉikkaṭṭi) — ice}

\label{lesson:TA-C184-panikkatti}



\begin{quote}
Before the new one: say the Tamil for a chutney, then the Tamil for a dry stir-fried vegetable dish.
\end{quote}

\subsection*{You'll want to know: \ta{பனிக்கட்டி}}
\textbf{\ta{பனிக்கட்டி}} — \emph{paṉikkaṭṭi} — \textquotedblleft{}ice\textquotedblright{}.

\begin{grammarlens}[title={What you've built}]
One more word for this chapter.
\end{grammarlens}

\subsection*{Your turn}
\begin{itemize}
\raggedright
  \item \emph{Say it:} \emph{paṉikkaṭṭi}
  \item \emph{Say it:} \emph{paṉikkaṭṭi}, once more
  \item \emph{Say it:} say \emph{paṉikkaṭṭi}
  \item \emph{Recall:} say the Tamil for a chutney, then the Tamil for a dry stir-fried vegetable dish, then say \emph{paṉikkaṭṭi} again
\end{itemize}

\begin{tcolorbox}[breakable,colback=teal!4,colframe=teal!35!black,title={Before you move on}]
What does \ta{பனிக்கட்டி} mean? (\textquotedblleft{}Ice\textquotedblright{}.) Say it once more.
\end{tcolorbox}
\section[\texorpdfstring{\ta{காலை} \ta{உணவு} (kālai uṇavu)}{காலை உணவு (kālai uṇavu)}]{\ta{காலை} \ta{உணவு} (kālai uṇavu) — breakfast}

\label{lesson:TA-C184-kalai-unavu}



\begin{quote}
Before the new one: say the Tamil for a dry stir-fried vegetable dish, then the Tamil for ice.
\end{quote}

\subsection*{You'll want to know: \ta{காலை} \ta{உணவு}}
\textbf{\ta{காலை} \ta{உணவு}} — \emph{kālai uṇavu} — \textquotedblleft{}breakfast\textquotedblright{}.

\begin{grammarlens}[title={What you've built}]
One more word for this chapter.
\end{grammarlens}

\subsection*{Your turn}
\begin{itemize}
\raggedright
  \item \emph{Say it:} \emph{kālai uṇavu}
  \item \emph{Say it:} \emph{kālai uṇavu}, once more
  \item \emph{Say it:} say \emph{kālai uṇavu}
  \item \emph{Recall:} say the Tamil for a dry stir-fried vegetable dish, then the Tamil for ice, then say \emph{kālai uṇavu} again
\end{itemize}

\begin{tcolorbox}[breakable,colback=teal!4,colframe=teal!35!black,title={Before you move on}]
What does \ta{காலை} \ta{உணவு} mean? (\textquotedblleft{}Breakfast\textquotedblright{}.) Say it once more.
\end{tcolorbox}
\section[\texorpdfstring{\ta{மதிய} \ta{உணவு} (matiya uṇavu)}{மதிய உணவு (matiya uṇavu)}]{\ta{மதிய} \ta{உணவு} (matiya uṇavu) — lunch}

\label{lesson:TA-C184-matiya-unavu}



\begin{quote}
Before the new one: say the Tamil for ice, then the Tamil for breakfast.
\end{quote}

\subsection*{You'll want to know: \ta{மதிய} \ta{உணவு}}
\textbf{\ta{மதிய} \ta{உணவு}} — \emph{matiya uṇavu} — \textquotedblleft{}lunch\textquotedblright{}.

\begin{grammarlens}[title={What you've built}]
One more word for this chapter.
\end{grammarlens}

\subsection*{Your turn}
\begin{itemize}
\raggedright
  \item \emph{Say it:} \emph{matiya uṇavu}
  \item \emph{Say it:} \emph{matiya uṇavu}, once more
  \item \emph{Say it:} say \emph{matiya uṇavu}
  \item \emph{Recall:} say the Tamil for ice, then the Tamil for breakfast, then say \emph{matiya uṇavu} again
\end{itemize}

\begin{tcolorbox}[breakable,colback=teal!4,colframe=teal!35!black,title={Before you move on}]
What does \ta{மதிய} \ta{உணவு} mean? (\textquotedblleft{}Lunch\textquotedblright{}.) Say it once more.
\end{tcolorbox}
\section[\texorpdfstring{\ta{இரவு} \ta{உணவு} (iravu uṇavu)}{இரவு உணவு (iravu uṇavu)}]{\ta{இரவு} \ta{உணவு} (iravu uṇavu) — dinner}

\label{lesson:TA-C184-iravu-unavu}



\begin{quote}
Before the new one: say the Tamil for breakfast, then the Tamil for lunch.
\end{quote}

\subsection*{You'll want to know: \ta{இரவு} \ta{உணவு}}
\textbf{\ta{இரவு} \ta{உணவு}} — \emph{iravu uṇavu} — \textquotedblleft{}dinner\textquotedblright{}.

\begin{grammarlens}[title={What you've built}]
One more word for this chapter.
\end{grammarlens}

\subsection*{Your turn}
\begin{itemize}
\raggedright
  \item \emph{Say it:} \emph{iravu uṇavu}
  \item \emph{Say it:} \emph{iravu uṇavu}, once more
  \item \emph{Say it:} say \emph{iravu uṇavu}
  \item \emph{Recall:} say the Tamil for breakfast, then the Tamil for lunch, then say \emph{iravu uṇavu} again
\end{itemize}

\begin{tcolorbox}[breakable,colback=teal!4,colframe=teal!35!black,title={Before you move on}]
What does \ta{இரவு} \ta{உணவு} mean? (\textquotedblleft{}Dinner\textquotedblright{}.) Say it once more.
\end{tcolorbox}
\section[\texorpdfstring{\ta{குறிப்பேடு} (kuṟippēṭu)}{குறிப்பேடு (kuṟippēṭu)}]{\ta{குறிப்பேடு} (kuṟippēṭu) — a notebook}

\label{lesson:TA-C184-kurippetu}



\begin{quote}
Before the new one: say the Tamil for lunch, then the Tamil for dinner.
\end{quote}

\subsection*{You'll want to know: \ta{குறிப்பேடு}}
\textbf{\ta{குறிப்பேடு}} — \emph{kuṟippēṭu} — \textquotedblleft{}a notebook\textquotedblright{}.

\begin{grammarlens}[title={What you've built}]
One more word for this chapter.
\end{grammarlens}

\subsection*{Your turn}
\begin{itemize}
\raggedright
  \item \emph{Say it:} \emph{kuṟippēṭu}
  \item \emph{Say it:} \emph{kuṟippēṭu}, once more
  \item \emph{Say it:} say \emph{kuṟippēṭu}
  \item \emph{Recall:} say the Tamil for lunch, then the Tamil for dinner, then say \emph{kuṟippēṭu} again
\end{itemize}

\begin{tcolorbox}[breakable,colback=teal!4,colframe=teal!35!black,title={Before you move on}]
What does \ta{குறிப்பேடு} mean? (\textquotedblleft{}A notebook\textquotedblright{}.) Say it once more.
\end{tcolorbox}
